% ============================================
% Predictability Analysis
% ============================================
\subsection*{H. Predictability Analysis}
\label{sec:suppl:predictability}

The Parental Dominance Index (PDI) and predictability analysis framework are defined in the main text Methods section. Here we provide additional details on the relationship between pairwise invasion outcomes and coalescence predictability.

% --------------------------------------------
% H.1 Dominant Species Identification
% --------------------------------------------
\subsubsection*{H.1 Dominant Species Identification}

For each coalescence event, we identified the dominant species in each parental community as the species with the highest relative abundance. When the dominant species from both parents were among the 12 isolates tested in pairwise invasion assays, we used their pairwise competition outcome to predict the coalescence direction.

% --------------------------------------------
% H.2 Predictability Across Conditions
% --------------------------------------------
\subsubsection*{H.2 Predictability Across Conditions}

Predictive power ($R^2$) between dominant-species competitive success and the Parental Dominance Index (PDI) was calculated separately for each nutrient condition:
\begin{itemize}
    \item \textbf{Nutr$-$}: $R^2 = 0.00$, no predictive power, consistent with weaker competitive interactions
    \item \textbf{Base}: $R^2 = 0.11$, weak predictability, suggesting emergent multi-species dynamics
    \item \textbf{Nutr$+$}: $R^2 = 0.49$, highest predictability, consistent with species-driven dynamics
\end{itemize}

This pattern supports the hypothesis that under strong-interaction conditions, coalescence outcomes are largely determined by competitive hierarchies among dominant species (\figref{fig:fig5}C).

